\chapter{The Verified Kernel Primitives in Depth}
\label{ch:kernel}

Chapters~\ref{ch:cmm} introduced the C{--} kernel at the layer level. This
chapter analyzes each primitive as a verified routine: its precondition,
algorithm, certificate, and failure mode.

\section{Memory Model: The Arena}
\label{sec:kern-arena}

All kernel allocation occurs in a fixed arena (\texttt{arena\_ptr},
\texttt{arena\_end}). There is no free list and no system allocator after
init. This has two consequences:
\begin{enumerate}
  \item \textbf{Bounded memory.} A computation cannot exhaust the heap
        unexpectedly; it either fits the arena or fails at allocation.
  \item \textbf{Determinism.} With no allocator metadata, two runs on the
        same input produce bit-identical layouts, which makes certificates
        reproducible.
\end{enumerate}

\section{Polynomial Addition}
\label{sec:kern-poly}

\paragraph{Precondition.} \texttt{p1}, \texttt{p2} point to valid
\texttt{[degree, c0\ensuremath{dots} cn]} blocks; the arena has room for a
$\max(d_1,d_2)+1$ coefficient result.
\paragraph{Algorithm.} Copy the larger degree; sum coefficient-wise, treating
missing high coefficients as zero.
\paragraph{Certificate.} The result satisfies
$\ensuremath{forall} i.\; \text{res}_i = p_{1,i} + p_{2,i}$, which the caller can re-check
against the input hashes stored in the WORM ledger.
\paragraph{Failure mode.} Arena exhaustion returns a pointer past
\texttt{arena\_end}; callers must bounds-check. (A hardened build would trap
here; the reference returns the advanced pointer for the caller to detect.)

\section{Matrix Multiplication}
\label{sec:kern-mat}

\paragraph{Precondition.} \texttt{m1} is $r_1\ensuremath{times} c_1$, \texttt{m2} is
$r_2\ensuremath{times} c_2$, $c_1 = r_2$.
\paragraph{Algorithm.} Triple loop with running sum; row-major indexing
$\text{data}[i\ensuremath{cdot} c + j + 2]$.
\paragraph{Safety boundary.} The explicit \texttt{if (c1 != r2) return 0}
enforces the dimension contract before any multiply. This is the single most
important correctness check in the numeric layer.
\paragraph{Certificate.} The result equals the definition
$(AB)_{ij} = \sum_k A_{ik}B_{kj}$; the FFI bridge verifies this against
Lean (\cref{ch:ffi}).

\section{Symbolic Differentiation}
\label{sec:kern-diff}

The primitive handles \texttt{TAG\_INT} and \texttt{TAG\_VAR} exactly and
returns the node for composite tags as a placeholder for the product-rule
dispatch. The design intent is that a later pass (or the Lean side) completes
the rule. The safety boundary here is tag validity: an unknown tag is not
dereferenced; the function returns the input unchanged rather than crashing.

\section{Algebraic Decision Procedure}
\label{sec:kern-dec}

\texttt{is\_equal\_verified} performs structural comparison and, in a full
build, would consult an SMT/SBV backend for algebraic equality (e.g., that
$x(x+1)/2 = (x^2+x)/2$). The reference stub uses pointer equality, which is
correct for canonical IR (since canonicalization makes semantic equality
syntactic) but would be unsound for non-canonical input. The boundary contract
therefore requires \emph{canonical} input --- a requirement enforced by
Layer~2.

\section{Normalization with Certificate}
\label{sec:kern-norm}

\texttt{normalize\_with\_cert} is the inverse-adjacent of canonicalization:
it expands and reorders, then writes a receipt seal. Because both Layer~2 and
Layer~5 normalize, the system has two independent normalizers that must agree
--- another cross-check.

\section{Summary of the Verified Surface}
\label{sec:kern-summary}

\begin{table}[H]
\centering
\caption{Kernel primitives and their safety boundaries.}
\label{tab:kern}
\begin{tabularx}{\textwidth}{lXl}
\toprule
\textbf{Primitive} & \textbf{Boundary enforced} & \textbf{Certificate} \\
\midrule
\texttt{arena\_init}      & fixed memory region          & memory bounds \\
\texttt{rewrite\_with\_cert} & certificate slot required & SHA-256 + Merkle \\
\texttt{poly\_add\_verified} & arena capacity           & coefficient law \\
\texttt{mat\_mul\_verified}  & dimension equality       & product law \\
\texttt{differentiate}       & tag validity             & derivative law \\
\texttt{is\_equal\_verified} & canonical input          & equality law \\
\texttt{normalize\_with\_cert} & expansion + seal       & receipt seal \\
\bottomrule
\end{tabularx}
\end{table}
This is a small but complete verified numeric kernel: every primitive names
its boundary and, where feasible, its certificate.
